\documentclass[letterpaper]{article}
\PassOptionsToPackage{unicode=true}{hyperref} % options for packages loaded elsewhere
\PassOptionsToPackage{hyphens}{url}
\usepackage{graphicx} % Required for inserting images


\def\tightlist{}


\graphicspath{ {../images/} }
\usepackage[margin=1in]{geometry}
\usepackage{tikz}
\usepackage[dvipsnames]{xcolor}
\usetikzlibrary{calc}




\usepackage[utf8]{inputenc}
\usepackage[T1]{fontenc}
\usepackage{lmodern}
\usepackage[english]{babel}


% Useful packages
\usepackage{hyperref}
\usepackage{amsmath, amssymb}
\usepackage{fancyhdr}
\usepackage{setspace}
\usepackage{tocloft}
\usepackage{titlesec}

% Code highlighting (Pandoc uses listings or minted, but listings is safer without shell escape)
\usepackage{listings}
\lstset{
  basicstyle=\ttfamily\small,
  breaklines=true,
  frame=single,
  backgroundcolor=\color[gray]{0.95}
}



%Modifiable Commands; defaults set to none
\newcommand{\titleDOC}{}
\newcommand{\authorDOC}{Daniel Sabalakov}
\newcommand{\dateDOC}{\today}
\newcommand{\classDOC}{}
% DEFAULT QUOTE IS BY ISSAC NEWTON. IF YOU DON'T WANT IT, DO: quote: @none
\newcommand{\quoteDOC}{``\textit{Truth is ever to be found in the simplicity, and not in the multiplicity and confusion of things}'' - Issac Newton}

% Header/Footer
\pagestyle{fancy}
\fancyhf{}
\rhead{\thepage}
\lhead{\leftmark}





% %Title Arguments
% \title{\TITLEOFDOCUMENT}
% \author{\AUTHOROFDOCUMENT}
% \date{\today}


%
% ANYTHING BELOW HERE CAN GET PROCEDURALLY GENERATED
%
% BEGINNING OF DOCUMENT
%
%
%
\begin{document}

%titlepage








\renewcommand{\classDOC}{DE Intro to Statistics I}

\renewcommand{\titleDOC}{Down the River}

\renewcommand{\authorDOC}{Daniel Sabalakov, Anthony Smith}





\begin{titlepage}
  \titleDOC
  \classDOC
  \dateDOC
  \authorDOC
\end{titlepage}

\begin{figure}[htbp]
	 \centering
	 \includegraphics[width=0.51\textwidth]{C:/Users/score/Downloads/imagelatex.jpg}
\end{figure}

\subsubsection{1. The authors use an analogy about the relationship
between a sample and a population. They introduce this analogy at the
beginning of the chapter and use it throughout the chapter. Identify the
analogy.}\label{the-authors-use-an-analogy-about-the-relationship-between-a-sample-and-a-population.-they-introduce-this-analogy-at-the-beginning-of-the-chapter-and-use-it-throughout-the-chapter.-identify-the-analogy.}

Vegetable soup; To decide whether or not it meets my standards, you can
just try a spoonful or two. This should be representative of the full
pot.

\subsubsection{2. ``Bias is the bane of sampling-the one thing above all
to avoid.'' BVD, p. 271. Define bias from a sampling perspective in one
sentence.}\label{bias-is-the-bane-of-sampling-the-one-thing-above-all-to-avoid.-bvd-p.-271.-define-bias-from-a-sampling-perspective-in-one-sentence.}

Bias is introduced by allowing your own ideas to influence sampling;
conclusions based on bias are unsalvageable and flawed.

\subsubsection{3. ``To avoid bias and make the sample as representative
as possible, you might be tempted to handpick the individuals included
in the sample with care and precision. The best strategy is to do
something quite different.'' BVD, p.
272.}\label{to-avoid-bias-and-make-the-sample-as-representative-as-possible-you-might-be-tempted-to-handpick-the-individuals-included-in-the-sample-with-care-and-precision.-the-best-strategy-is-to-do-something-quite-different.-bvd-p.-272.}

 \textbf{ randomly. } 

\subsubsection{4. The authors wrote about five legitimate sampling
methods. Identify as many as you
can.}\label{the-authors-wrote-about-five-legitimate-sampling-methods.-identify-as-many-as-you-can.}

\begin{itemize}
\tightlist
\item
  Truly random (simple)

  \begin{itemize}
  \tightlist
  \item
    Selecting people truly randomly
  \end{itemize}
\item
  Stratified

  \begin{itemize}
  \tightlist
  \item
    Identifying homogenous groups and then sampling from them
  \end{itemize}
\item
  Cluster

  \begin{itemize}
  \tightlist
  \item
    Splitting population into similar parts and them sampling them
  \end{itemize}
\item
  Systematic

  \begin{itemize}
  \tightlist
  \item
    Sampling population using systematic methods
  \end{itemize}
\end{itemize}

\subsubsection{5. Choose one of the sampling methods from your answer to
question four and describe
it.}\label{choose-one-of-the-sampling-methods-from-your-answer-to-question-four-and-describe-it.}

Cluster: When splitting a population into multiple easy-to-manage parts
and then sampling a few random points from each cluster.

\subsubsection{6. If samples are drawn at random, would we expect them
to generally differ one from another? Briefly
explain.}\label{if-samples-are-drawn-at-random-would-we-expect-them-to-generally-differ-one-from-another-briefly-explain.}

Yes; truly random samples will differ from each other because each
combination has the same chance of being drawn.

\subsubsection{7. ``Bad sample designs yield worthless data. Many of the
most convenient forms of sampling can be seriously biased. And there is
no way to correct for the bias from a bad sample.'' BVD, p. 282. The
authors wrote about four major forms of bias. Identify as many as you
can.}\label{bad-sample-designs-yield-worthless-data.-many-of-the-most-convenient-forms-of-sampling-can-be-seriously-biased.-and-there-is-no-way-to-correct-for-the-bias-from-a-bad-sample.-bvd-p.-282.-the-authors-wrote-about-four-major-forms-of-bias.-identify-as-many-as-you-can.}

\begin{itemize}
\tightlist
\item
  Volunteer Sample Bias
\item
  Convenience sampling
\item
  Bad Sample Frame
\item
  Undercoverage
\end{itemize}

\subsubsection{8. Choose one of the forms of bias and write about
it.}\label{choose-one-of-the-forms-of-bias-and-write-about-it.}

Undercoverage is where some portion of the population is not sampled at
all. This can arise from any reason.



\end{document}
